\documentclass[11pt]{letter}
\usepackage[utf8]{inputenc}
\usepackage[T1]{fontenc}
\usepackage[english]{babel}
\usepackage[a4paper,margin=1in]{geometry}

\signature{Bartosz Lewandowski, Ph.D.\\
Department of Clinical Engineering, Academy of Silesia in Katowice\\
bartosz.lewandowski@akademiaslaska.pl}
\address{}

\begin{document}
\begin{letter}{The Editors\\ \textit{Applied Acoustics}}

\opening{Dear Editors,}

We are pleased to submit the manuscript entitled ``Assessing the Similarity of Acoustic Signals Using a Vector Representation of Sound Features -- From Acoustic Features to Perceptual Difficulty: How to Compare Two Sounds?'' for consideration as a Full Research Article in \textit{Applied Acoustics}.

Auditory tests that present several sounds simultaneously, for example tests of auditory perception and cognitive reserve, assume that task difficulty reflects the ability being measured. If two stimuli are acoustically very close, however, the test may measure the similarity of the sound material rather than the listener. In this work we describe each of 24 animal sounds by a 16-dimensional vector of acoustic features (spectral centroid, fundamental frequency, peak-to-average ratio and 13 mel-frequency cepstral coefficients), use the Euclidean distance between min-max normalized vectors as an objective measure of acoustic proximity, and compare it with the behavioural results of the Memory Load Test (MLT) study. The stimulus pair with the smallest distance was identified correctly in only 20.5\% of trials, whereas all more distant pairs reached 44.4--92.9\%; the relationship was not monotonic, which we report and discuss openly.

The manuscript fits the scope of the journal in the areas of acoustic signal processing and psychoacoustics. Its practical contribution is a simple, reproducible and auditable procedure for screening stimulus sets for overly similar pairs before they are used in audiological and cognitive testing.

We confirm that this work has not been published previously, is not under consideration for publication elsewhere, and that its publication has been approved by all authors.

\closing{Yours sincerely,}

\end{letter}
\end{document}
